% ============================================================
%  TEMA 2 — APARTADO 3: FORMAS DE REDISEÑO DE PUESTOS
%
%  RECORDATORIO AL EQUIPO:
%  Las principales formas de rediseño de puestos son:
%
%    1. ROTACIÓN DE PUESTOS (job rotation): los empleados rotan entre puestos.
%       Ventajas: polivalencia, motivación, reducción de la monotonía.
%       Inconvenientes: costes de formación, pérdida de especialización.
%
%    2. AMPLIACIÓN DE TAREAS (job enlargement / ensanchamiento horizontal):
%       se añaden tareas del mismo nivel de complejidad.
%       Ventajas: reduce la monotonía, mayor variedad.
%       Inconvenientes: puede aumentar la carga sin añadir motivación real.
%
%    3. ENRIQUECIMIENTO DE PUESTOS (job enrichment / Hackman & Oldham):
%       se añaden responsabilidades de mayor complejidad (carga vertical).
%       Las cinco características centrales del modelo:
%         - Variedad de habilidades
%         - Identidad de la tarea
%         - Significado de la tarea
%         - Autonomía
%         - Retroalimentación
%       Ventajas: mayor motivación intrínseca, satisfacción, compromiso.
%       Inconvenientes: requiere empleados que quieran / puedan asumir más.
%
%    4. EQUIPOS DE TRABAJO AUTOGESTIONADOS (self-managed teams):
%       grupos con alta autonomía para organizar su propio trabajo.
%
%    5. TELETRABAJO / TRABAJO FLEXIBLE (flexibilización del puesto):
%       modificación del lugar y/o horario de trabajo.
%
%    6. JOB CRAFTING: el empleado adapta el puesto a sus fortalezas
%       de forma proactiva (tendencia emergente).
%
%  PARA CADA FORMA INDICAD:
%  - ¿La aplica CaixaBank? ¿En qué puestos / áreas?
%  - Ventajas e inconvenientes concretos observados.
%  - ¿Está la empresa satisfecha?
%  - ¿Qué formas serían interesantes aplicar y por qué?
% ============================================================

\section{Formas de rediseño de puestos en CaixaBank}

\subsection{Formas de rediseño aplicadas actualmente}

% ---- TABLA RESUMEN ----
\begin{table}[H]
  \centering
  \small
  \caption{Aplicación de formas de rediseño de puestos en CaixaBank}
  \begin{tabularx}{\textwidth}{p{0.28\textwidth} p{0.12\textwidth} X}
    \toprule
    \textbf{Forma de rediseño}      & \textbf{¿Se aplica?} & \textbf{Descripción y observaciones} \\
    \midrule
    Rotación de puestos             & [Sí/No/Parcial] & \\
    Ampliación de tareas            & [Sí/No/Parcial] & \\
    Enriquecimiento de puestos      & [Sí/No/Parcial] & \\
    Equipos autogestionados         & [Sí/No/Parcial] & \\
    Teletrabajo / trabajo flexible  & [Sí/No/Parcial] & \\
    Job crafting                    & [Sí/No/Parcial] & \\
    \bottomrule
  \end{tabularx}
\end{table}

% ---- ROTACIÓN ----
\subsubsection{Rotación de puestos (\textit{job rotation})}

% TODO: ¿Existe un programa formal de rotación en CaixaBank?
% Es habitual en banca para desarrollar directivos y gestores con visión global.
% ¿Se rota entre oficinas, entre departamentos, internacionalmente (BPI)?

[Descripción a completar]

\textbf{Ventajas observadas en CaixaBank:}
\begin{itemize}
  \item [...]
\end{itemize}

\textbf{Inconvenientes observados:}
\begin{itemize}
  \item [...]
\end{itemize}

% ---- AMPLIACIÓN ----
\subsubsection{Ampliación de tareas (\textit{job enlargement})}

% TODO: Con la digitalización bancaria, los asesores de oficina han asumido
% tareas que antes hacían otros departamentos (apertura online, soporte digital...).
% ¿Esto genera sobrecarga? ¿Hay compensación retributiva?

[Descripción a completar]

% ---- ENRIQUECIMIENTO ----
\subsubsection{Enriquecimiento de puestos (\textit{job enrichment})}

% TODO: ¿Se da autonomía a los asesores para tomar decisiones sobre productos?
% ¿Tienen feedback claro sobre su impacto en los resultados?
% El modelo de Hackman y Oldham es muy útil para evaluar si el diseño actual
% de los puestos genera motivación intrínseca.

[Descripción a completar]

% ---- TELETRABAJO ----
\subsubsection{Teletrabajo y trabajo flexible}

% TODO: Tras la pandemia COVID-19, CaixaBank implementó políticas de teletrabajo.
% ¿Sigue vigente? ¿Para qué puestos? ¿Con qué porcentaje de jornada?
% ¿Ha mejorado la satisfacción y productividad? ¿Ha generado problemas
% de desconexión del equipo o de la cultura organizacional?

[Descripción a completar]

\subsection{Formas de rediseño recomendadas para CaixaBank}

% TODO: Proponed al menos 2-3 formas que no se apliquen o que se apliquen
% de forma mejorable. Justificad por qué serían beneficiosas.
% Ejemplos posibles:
%   - Impulsar el job crafting en puestos muy regulados para mejorar motivación.
%   - Aplicar equipos autogestionados en áreas de innovación/tecnología.
%   - Revisar el job enlargement para evitar sobrecarga sin contraprestación.

\begin{notaequipo}
  Preguntad en la entrevista:
  \begin{enumerate}
    \item ¿Existe algún programa formal de rotación? ¿Tiene criterios claros?
    \item ¿Los empleados de oficina han asumido más tareas en los últimos años?
          ¿Se ha revisado su carga de trabajo y retribución?
    \item ¿Tienen autonomía para adaptar su forma de trabajo (\textit{job crafting})?
    \item ¿Cuál es la política actual de teletrabajo? ¿Están satisfechos con ella?
    \item ¿Se mide el impacto de los rediseños en productividad y satisfacción?
  \end{enumerate}
\end{notaequipo}
